% Part: methods
% Chapter: proofs
% Section: reading-proofs

\documentclass[../../../include/open-logic-section]{subfiles}

\begin{document}

\olfileid[hi]{mth}{prf}{rea}
\olsection{प्रमाण पढ़ना}

पाठ्यपुस्तकों और लेखों में मिलने वाले प्रमाण बहुत कम उन सभी विवरणों को देते
हैं जिन्हें अब तक हमने अपने उदाहरणों में शामिल किया है। लेखक प्रायः यह
ध्यान नहीं दिलाते कि वे कब स्थितियों में भेद करते, कब अप्रत्यक्ष प्रमाण देते,
या यह उल्लेख नहीं करते कि परिभाषा उपयोग कर रहे हैं। अतः पाठ्यपुस्तक में
प्रमाण पढ़ते समय उसे समझने के लिए अक्सर वे विवरण स्वयं भरने होंगे। यह करना
प्रमाण में आवश्यक विभिन्न चालों का अभ्यास भी है। एक उदाहरण देखें।

\begin{prop}[अवशोषण]
सभी समुच्चयों $A$, $B$ के लिए,
\[
A \cap (A \cup B) = A
\]
\end{prop}

\begin{proof}
यदि $z\in A\cap(A\cup B)$, तो $z\in A$, अतः
$A\cap(A\cup B)\subseteq A$। अब मान लें $z\in A$। तब
$z\in A\cup B$ भी, और इसलिए $z\in A\cap(A\cup B)$ भी।
\end{proof}

अवशोषण नियम का पिछला प्रमाण बहुत संक्षिप्त है। प्रयुक्त परिभाषाओं का कोई
उल्लेख नहीं, सिद्ध करने से पहले ``हमें यह सिद्ध करना है'' नहीं, इत्यादि। इसे
खोलें। सिद्ध प्रतिज्ञप्ति किन्हीं समुच्चयों $A$ और $B$ के बारे में सामान्य
दावा है और प्रमाण में $A$ या $B$ का उल्लेख मनमाने समुच्चयों के चर हैं। प्रमाण
जो सामान्य दावे स्थापित करता है वे समुच्चयों की समानता सिद्ध करने के लिए
आवश्यक हैं, अर्थात् समानता के बाएँ पक्ष का प्रत्येक !!{element}, दाएँ पक्ष
का !!a{element} है और इसके विपरीत।

\begin{quote}
``यदि $z\in A\cap(A\cup B)$, तो $z\in A$, अतः
$A\cap(A\cup B)\subseteq A$।''
\end{quote}

यह समानता के प्रमाण का पहला आधा है: यह स्थापित करता है कि यदि कोई मनमाना
$z$ बाएँ पक्ष का !!a{element} है, तो वह दाएँ पक्ष का भी !!a{element} है,
अर्थात् $A\cap(A\cup B)\subseteq A$। मान लें
$z\in A\cap(A\cup B)$। चूँकि $z$ दो समुच्चयों के सर्वनिष्ठ का
!!{element} तभी है जब दोनों समुच्चयों का !!{element} हो, हम निष्कर्षित कर
सकते हैं कि $z\in A$ और $z\in A\cup B$ भी। विशेषतः $z\in A$, जिसे दिखाना
था। चूँकि पहले आधे के लिए इतना ही करना है, हम जानते हैं कि शेष प्रमाण दूसरे
आधे, अर्थात् $A\subseteq A\cap(A\cup B)$, का प्रमाण होना चाहिए।

\begin{quote}
``अब मान लें $z\in A$। तब $z\in A\cup B$ भी, और इसलिए
$z\in A\cap(A\cup B)$ भी।''
\end{quote}

हम $z\in A$ मानकर आरंभ करते हैं, क्योंकि दिखा रहे हैं कि किसी भी $z$ के
लिए यदि $z\in A$, तो $z\in A\cap(A\cup B)$। $z\in A\cap(A\cup B)$ दिखाने
के लिए (``$\cap$'' की परिभाषा से) दिखाना होगा कि (i) $z\in A$ और (ii)
$z\in A\cup B$ भी। यहाँ (i) केवल हमारी मान्यता है, अतः आगे कुछ सिद्ध नहीं
करना और इसीलिए प्रमाण उसका फिर उल्लेख नहीं करता। (ii) के लिए याद करें कि
$z$ समुच्चयों के सम्मिलन का !!a{element} तभी है जब उन समुच्चयों में
कम-से-कम एक का !!{element} हो। चूँकि $z\in A$ और $A\cup B$, $A$ तथा $B$
का सम्मिलन है, यहाँ ऐसा ही है। अतः $z\in A\cup B$। हमने (i) $z\in A$ और
(ii) $z\in A\cup B$ दोनों दिखाए, फलतः ``$\cap$'' की परिभाषा से
$z\in A\cap(A\cup B)$। प्रमाण उन परिभाषाओं का उल्लेख नहीं करता; माना जाता
है कि पाठक उन्हें आत्मसात कर चुका है। यदि नहीं, तो लौटकर स्वयं को याद दिलाना
होगा कि वे क्या हैं। फिर यह भी पहचानना होगा कि $z\in A$ से $z\in A\cup B$
क्यों निष्पन्न होता है, और $z\in A$ तथा $z\in A\cup B$ से
$z\in A\cap(A\cup B)$ क्यों।

ऊपर के प्रमाण का एक अन्य रूप यहाँ है, जिसमें सब कुछ स्पष्ट कर दिया गया है:
\begin{proof}{}
[समुच्चयों के लिए $=$ की परिभाषा से $A \cap (A \cup B) = A$ सिद्ध करने
  के लिए हमें दिखाना है कि (a) $A \cap (A \cup B) \subseteq A$ और (b)
  $A \cap (A \cup B) \subseteq A$। (a): $\subseteq$ की परिभाषा से हमें
  दिखाना है कि यदि $z \in A \cap (A \cup B)$, तो $z \in A$।] यदि $z
\in A \cap (A \cup B)$, तो $z \in A$ [क्योंकि $\cap$ की परिभाषा से
  $z \in A \cap (A \cup B)$ तभी और केवल तभी जब $z \in A$ और $z \in
  A \cup B$], अतः $A \cap (A \cup B) \subseteq A$। [(b): $\subseteq$
  की परिभाषा से हमें दिखाना है कि यदि $z \in A$, तो $z \in A \cap (A
  \cup B)$।] अब मान लें [(1)] $z \in A$। तब [(2)] $z \in A \cup B$
भी [क्योंकि (1) से $z \in A$ या $z \in B$, जिसका $\cup$ की परिभाषा से
  अर्थ है $z \in A \cup B$], और इसलिए $z \in A \cap (A \cup B)$ भी
[क्योंकि $\cap$ की परिभाषा अपेक्षा करती है कि $z \in A$, अर्थात् (1),
  और $z \in A \cup B)$, अर्थात् (2)]।
\end{proof}

\begin{prob}
$A \cup (A \cap B) = A$ के निम्नलिखित प्रमाण को विस्तृत कीजिए। इसमें
प्रयुक्त सभी अनुमान प्रतिरूपों का उल्लेख कीजिए, बताइए कि प्रत्येक चरण पहले
स्थापित मान्यताओं या दावों से क्यों निष्पन्न होता है, और कहाँ हमें किन
परिभाषाओं का सहारा लेना पड़ता है।
\begin{proof}
  यदि $z \in A \cup (A \cap B)$, तो $z \in A$ या $z \in A \cap B$।
  यदि $z \in A \cap B$, तो $z \in A$। प्रत्येक $z \in A$, $A \cup (A
  \cap B)$ में भी है।
\end{proof}
\end{prob}

\end{document}
